% Redrawn 2026-10-15; main labels 8 pt, notes 7.5 pt; actual fonts embedded in PDF.
% export-reference-width-mm: 169
\begin{tikzpicture}[font=\figurefont\fontsize{8}{10}\selectfont,
  box/.style={line width=.6pt,rounded corners=1.8pt,draw=FigureStroke,fill=FigureBlueFill,align=center,
    minimum width=19mm,minimum height=10mm},
  arr/.style={-{Latex[length=1.8mm]},draw=FigureStroke,line width=.6pt}]
  \node[anchor=west,font=\figurefont\mdseries] at (-1,3) {(a) 误差落在关键岔路};
  \node[box,fill=FigureBlueFill] (start) at (0,1) {起点};
  \node[box,fill=FigureRedFill] (fork) at (2.7,1) {岔路\\训练抽样权重低};
  \node[box,fill=FigureBlueFill] (left) at (5.8,2) {真实1\\估计0};
  \node[box] (right) at (5.8,0) {真实0\\估计1};
  \draw[arr] (start) -- (fork);
  \draw[arr,dashed] (fork) -- (left);
  \draw[arr] (fork) -- (right);
  \node[align=center] at (2.8,-1) {平均平方误差$10^{-3}$\\部署损失仍可为1};
  \begin{scope}[shift={(9.5,0)}]
    \node[anchor=west,font=\figurefont\mdseries] at (-1,3) {(b) 估计与访问的反馈};
    \node[box,fill=FigureBlueFill] (data) at (0,1.8) {样本与覆盖};
    \node[box,fill=FigureGreenFill] (q) at (3.5,1.8) {$q$估计};
    \node[box,fill=FigureBlueFill] (policy) at (3.5,-.3) {动作与访问};
    \node[box] (class) at (0,-.3) {表示与优化};
    \draw[arr] (data) -- (q);
    \draw[arr] (class) -- (q);
    \draw[arr] (q) -- (policy);
    \draw[arr,dashed] (policy.south) -- (3.5,-1.4)
      -- (-1.4,-1.4) -- (-1.4,1.8) -- (data.west);
    \node[fill=white] at (1,-1.4) {分布改变};
  \end{scope}
\end{tikzpicture}
